% r03_hyb_A0: proof of \Cref{rot:prop:BD}, part IV: the A=0 input (the hybrid kernels on Lambda_X and on periodic cells) and the pointwise and Hoelder bounds at a regular background on a cell.

\paragraph{The $A=0$ input and the bounds on a cell.}

\begin{lemma}[$A=0$ bounds for the hybrid flow]\label{rot:hyb:lem:A0}
On $\Lambda_X$, with $S^X_j$ the slice kernels of \Cref{rot:def:match} ($G^X_r(0)=\sum_{j<r}S^X_j$ exactly):
\begin{enumerate}[label=\textup{(\alph*)},leftmargin=*]
\item $|S^X_j(x,y)|\le C\ell_j^{-1}e^{-\delta|x-y|/\ell_j}$, $|\nabla_xS^X_j(x,y)|\le C\ell_j^{-2}e^{-\delta|x-y|/\ell_j}$, and for $0<\alpha<1$
\[
|\nabla_xS^X_j(x,y)-\nabla_xS^X_j(x',y)|\le C_\alpha\log^{3\alpha}(2+\ell_j/\eta')\,\ell_j^{-2-\alpha}|x-x'|^\alpha e^{-\delta\min(|x-y|,|x'-y|)/\ell_j}
\]
(the logarithm is the slice's own, $\le\log^{3\alpha}(1/\eta')$);
\item consequently $\|G^X_r(0)f\|_\infty+\|\nabla G^X_r(0)f\|_\infty\le C\|f\|_\infty$ with $C$ uniform in $r$, $[\nabla G^X_r(0)f]_\alpha\le C_\alpha\log^{3\alpha}(1/\eta')\|f\|_\infty$, and $\|G^X_r(0)f\|_q\le C\|f\|_p$ for $0\le\frac1p-\frac1q<\frac23$, with the exponential localisation $e^{-\delta_0\dist(x,\supp f)}$ (\cite[(2.16)--(2.17)]{Ba4});
\item on a periodic cell $\square=\Lambda_X/P$ fitted by both lattices, of side $\ge2M$, and for a constant field $B_0$,
\[
G^X_r(\square,B_0)(x,y)=e^{\mathrm ie_rB_0\cdot(x-y)}\sum_{\gamma\in P}e^{\mathrm ie_rB_0\cdot\gamma}G^X_r(0)(x,y+\gamma)
\]
(twisted periodisation), and (b) holds on $\square$ for $G^X_r(\square,B_0)$ with constants independent of $B_0$ and $M$, the H\"older constant carrying the same $\log^{3\alpha}(1/\eta')$.
\end{enumerate}
\end{lemma}
\begin{proof}
(a) In the unit scale of level $j$, $S^X_j=H^X_jC^X_jH^{X*}_j$ (\Cref{rot:def:match}). The decay $|H^X_j(x,z)|\le Ce^{-\delta|x-z|}$ and the gradient bound $|\nabla H^X_j|\le C_\nabla$ are \Cref{rot:lem:H,rot:lem:grad} (the decay is the Combes--Thomas decay of the coercive coarse operators, \cite{CT73}); the decay of the coarse kernel $C^X_j=(A^X_j+bQ^*Q)^{-1}$ is the one used in the proof of \Cref{rot:lem:slice} (its symbol is analytic and bounded on a strip because $A^X_j+bQ^*Q\ge c$). \emph{The H\"older derivative.} The second difference of $H^X$ in $x$ multiplies the alias-$l$ summand of $M^X$ in the representation of \Cref{rot:lem:H} by $d(q+l)^2$ with $|d(Q)|\le|Q|$; since $\Delta'_X(Q)\ge c|Q|^2$ on each alias cell and $\sum_lu_X(q+l)^2=1$, the alias sum is $\le C\sum_l|u_X(q+l)|\le C$, and integrating against $|u_{r-1}|$ over $V_\eta$ costs $J_0\le C\log^3(1/\eta)$ (\Cref{rot:lem:CK0}). So $|\nabla\nabla H^X|\le\Lambda_\eta\defeq C\log^3(1/\eta)$ and, by interpolation with $|\nabla H^X|\le C_\nabla$,
\[
|\nabla H^X(x,z)-\nabla H^X(x',z)|\le\min\{2C_\nabla,\Lambda_\eta|x-x'|\}\le(2C_\nabla)^{1-\alpha}\Lambda_\eta^{\alpha}|x-x'|^\alpha\le C_\alpha\log^{3\alpha}(1/\eta)|x-x'|^\alpha
\]
for every $\alpha<1$; the geometric mean with the decay gives the localised form. Composing with $C^X_j$ and $H^{X*}_j$ (exponentially decaying kernels) and rescaling by \cite[(2.20)]{K1} (scalar kernels scale as $\ell^{-1}$ in $d=3$, derivatives as $\ell^{-2}$) gives (a); at level $j$ the logarithm is $\log(1/\eta_j)$ with $\eta_j=5^{1-j}$, i.e.\ $\le\log(2+\ell_j/\eta')$. \emph{Summation.} $\|S^X_jf\|_\infty\le C\ell_j^{-1}\ell_j^3\|f\|_\infty$, $\|\nabla S^X_jf\|_\infty\le C\ell_j\|f\|_\infty$, $[\nabla S^X_jf]_\alpha\le C_\alpha\log^{3\alpha}(2+\ell_j/\eta')\,\ell_j^{1-\alpha}\|f\|_\infty$; summing over $j<r$ converges for $\alpha<1$, the H\"older term to $\le C_\alpha\sum_{j<r}5^{-(r-j)(1-\alpha)}(1+j\ln5)^{3\alpha}\le C_\alpha\log^{3\alpha}(1/\eta')$ (this is \cite[(2.38)--(2.39)]{Ba4} with the slice kernels in place of $G_jQ_j^*C^{(j)}Q_jG_j$). The $L^p\to L^q$ bounds follow from the kernel bound $|G^X_r(0)(x,y)|\le C\min(\eta'^{-1},|x-y|^{-1})e^{-\delta|x-y|}$ by Young's inequality and the Riesz--Thorin theorem (\cite[p.~583]{Ba4}). (c) Both sides are the inverse of $-\Delta_{B_0}+m_r^2+a_rP_X(B_0)$ on $P$-periodic functions: $P\subset\Gamma_r$ makes the hybrid block system $P$-periodic, and $e^{\mathrm ie_rB_0\cdot x}$ conjugates the operator to the $A=0$ operator on functions with the twisted periodicity $\psi(x+\gamma)=e^{-\mathrm ie_rB_0\cdot\gamma}\psi(x)$, whose kernel is the twisted image sum (\cite[(B.3)]{K2} with phases). The images with $\gamma\ne0$ contribute $\le Ce^{-\delta_0(2M-\sqrt3)}$ to every bound in (b), uniformly in $B_0$.
\end{proof}

\begin{lemma}[{\cite[Lemma~2.2]{Ba4}, \cite[Lemma~5.8]{K2}, for the hybrid flow}]\label{rot:hyb:lem:HOL}
Let $\square$ be a periodic cell fitted by both lattices of side $2M$, and $\tilde A=A_0+A'$ regular on $\square$, $A_0$ constant, with $|A'|_\infty,|\partial A'|_\infty\le c'e_r^{\beta-1}$ and $A'$ supported away from the boundary of the fundamental region. Then for $e_r^\beta\le c_0(a,M)$, with $c_0,c_1,c_2$ uniform in $r$,
\[
\|G^X_r(\square,\tilde A)f\|_{1,0}\le c_1\|f\|_\infty,\qquad\|G^X_r(\square,\tilde A)f\|_{1,\alpha}\le c_1C_\alpha\log^{3\alpha}(1/\eta')\|f\|_\infty,
\]
\[
\|G^X_r(\square,\tilde A)f\|_q,\ \|D_{\tilde A}G^X_rf\|_q,\ \|G^X_rD_{\tilde A}^*f\|_q\le c_2\|f\|_p
\]
for $1\le p,q\le\infty$ with $\frac1p-\frac1q<\frac23$, where $\|\cdot\|_{1,\alpha}$ is the covariant H\"older norm \cite[(2.14)]{Ba4} with the field $\tilde A$ and $\|\cdot\|_{1,0}$ the sup norm plus the sup of the covariant derivative.
\end{lemma}
\begin{proof}
\cite[(2.24)--(2.26), (2.31)--(2.33)]{Ba4}: $G_r(\square,\tilde A)=\sum_{n\ge0}G_r(\square,A_0)[V_rG_r(\square,A_0)]^n$ with $V_r=V_r(A',A_0)$ the vertex operator of \cite[(3.16)]{Ba1} ($=$ \cite[(3.53)]{K1}); $V_r$ is a first-order differential operator with coefficients $O(e_r|A'|,e_r|\partial A'|)$ (\cite[(2.32)]{Ba4}, using that $A'$ vanishes near the boundary of the fundamental region) plus the pin part $a_r[Q(A_0)^*F_{2,r}(A',A_0)+F_{2,r}^*Q(A_0)+F_{2,r}^*F_{2,r}]$ with $(F_{2,r}(A',A_0)\phi)(y)=\sum_{x\in B(y)}\eta'^3(U(A'(\Gamma_{y,x}))-1)U(A_0(\Gamma_{y,x}))^{-1}\phi(x)$, $\|F_{2,r}\|_{\infty\to\infty}\le\sup|U(A'(\Gamma))-1|\le2.1\,e_r|A'|_\infty$ by \Cref{rot:hyb:lem:len}(i). Hence $\|V_rg\|_\infty\le C(M)e_r^\beta\|g\|_{1,0}$. By \Cref{rot:hyb:lem:A0}(b),(c), $\|G_r(\square,A_0)g\|_{1,0}\le c\|g\|_\infty$ with $c$ uniform in $r$ and $A_0$, and $\|G_r(\square,A_0)g\|_{1,\alpha}\le cC_\alpha\log^{3\alpha}(1/\eta')\|g\|_\infty$ (the constant background is a twisted periodisation of the $A=0$ kernel, and the H\"older norm (2.14) with field $A_0$ is the plain H\"older norm after the local gauge $e^{\mathrm ie_rA_0\cdot x}$ on each unit cube). \emph{Order of summation.} Put $g_0\defeq f$, $g_n\defeq V_rG_r(\square,A_0)g_{n-1}$. Then $\|g_n\|_\infty\le C(M)e_r^\beta\|G_r(\square,A_0)g_{n-1}\|_{1,0}\le(C(M)c\,e_r^\beta)^n\|f\|_\infty$, so for $C(M)c\,e_r^\beta\le\frac12$, a condition uniform in $r$, the sum $g\defeq\sum_ng_n$ converges in the sup norm with $\|g\|_\infty\le2\|f\|_\infty$, and $G_r(\square,\tilde A)f=\sum_nG_r(\square,A_0)g_n=G_r(\square,A_0)g$ (the cell is finite, every operator is bounded, and the sums may be exchanged). Hence $\|G_r(\square,\tilde A)f\|_{1,0}\le2c\|f\|_\infty$ and $\|G_r(\square,\tilde A)f\|_{1,\alpha}\le2cC_\alpha\log^{3\alpha}(1/\eta')\|f\|_\infty$: the H\"older constant enters once, at the end, and never the smallness condition. (Summing in $\|\cdot\|_{1,\alpha}$, as \cite[(2.33)]{Ba4} may with its uniform constant, would turn the condition into $e_r^\beta\log^{3\alpha}(1/\eta')\le c_0$, which depends on $r$.) The norms (2.14) with the fields $A_0$ and $\tilde A$ differ on the cell by at most $Ce_r^\beta\|\cdot\|_{1,0}$ ($D_{\tilde A}=U(A')D_{A_0}+F_{1,r}(A')$ and $|U(A'(\Gamma_{y,x}))-1|\le Ce_r|A'|_\infty|x-y|$ for the bounded-length paths of (2.14)), which changes $c_1$ by a factor $1+O(e_r^\beta)$. The $L^p\to L^q$ bounds are summed in the same order, with $D_{\tilde A}=U(A')D_{A_0}+F_{1,r}(A')$.
\end{proof}

\begin{remark}[the H\"older constant]\label{rot:hyb:rem:Hol}
In \Cref{rot:prop:BD}(a) the H\"older clause \cite[(3.8)]{K1} of $G^X_r(A)$, and the H\"older clauses of $H^X(A)$ and $\Gamma^X(A)$, hold with the constant $C_\alpha\log^{3\alpha}(1/\eta')$, $0<\alpha<1$ (\Cref{rot:hyb:lem:A0}, \Cref{rot:hyb:lem:HOL}, and the random walk below); on slice $j$ the constant is $C_\alpha\log^{3\alpha}(2+\ell_j/\eta')$ (\Cref{rot:hyb:lem:A0}(a)). The sup, decay and one-derivative bounds (3.6)--(3.7) and the $\|\cdot\|_{1,0}$ bounds have constants uniform in $r$, and so does every smallness condition. \cite[(3.8)]{K1} asks for some $0<\alpha<1$ with some constant (\cite[p.~656]{K1}). Every use of the clause in this section is either as a gain, where $\eta'^\alpha\log^{3\alpha}(1/\eta')\le C_{\alpha,\alpha'}\eta'^{\alpha'}$ and, slice by slice, $(\eta'/\ell_j)^\alpha\log^{3\alpha}(2+\ell_j/\eta')\le C_{\alpha,\alpha'}(\eta'/\ell_j)^{\alpha'}$ for $\alpha'<\alpha$ (\Cref{rot:thm:H1S}, whose exponent $\theta_1$ is fixed by $\alpha$ alone), or inside a per-flow bound that already allows a factor $(1+k)^p$ (\Cref{rot:cor:OF}), where $\log^{3\alpha}(1/\eta')=(r\ln5)^{3\alpha}\le C(1+k)^{3\alpha}$.
\end{remark}
